\section{Related Works}
\label{sec:related}

\begin{table}[t]
\centering
\caption{Comparison of radar rendering methods. Signal: ADC\,=\,raw FMCW, CIR\,=\,channel impulse response, IF\,=\,intermediate-frequency phasors, RA\,=\,range--azimuth image, RDA\,=\,range--Doppler--azimuth cube. Complex\,=\,outputs complex-valued signals (not power-only). Repr.\,=\,scene representation; Mat.\,=\,physics material model; Per-pt\,=\,per-point/per-vertex materials; NVS\,=\,novel view synthesis demonstrated; Real\,=\,validated on real data. Time\,=\,per-scene wall-clock training time on a single RTX~4090 over the 8 training viewpoints used in our novel-view experiments (Sec.~\ref{sec:experiments}); blank entries are forward-only or single-viewpoint methods.}
\label{tab:comparison}
\small
\setlength{\tabcolsep}{3pt}
\renewcommand{\arraystretch}{1.1}
\resizebox{\linewidth}{!}{
\begin{tabular}{lcccccccccc}
\toprule
\textbf{Method} & \textbf{Repr.} & \textbf{Signal} & \textbf{Complex} & \textbf{Diff.} & \textbf{Mat.} & \textbf{Per-pt} & \textbf{BSDF} & \textbf{NVS} & \textbf{Real} & \textbf{Time} \\
\midrule
Sionna RT~\cite{sionna}              & Mesh     & CIR & \checkmark & Partial    & --         & --         & --         & --         & --         & --       \\
SFCW Inv.~\cite{hofmann2025inverse}  & Mesh     & IF  & \checkmark & Partial    & \checkmark & \checkmark & --         & --         & \checkmark & --       \\
InverTwin~\cite{chen2025invertwin}   & Mesh     & CIR & \checkmark & \checkmark & --         & --         & --         & --         & \checkmark & --       \\
mmIR~\cite{mmIR2026}                 & Mesh     & ADC & \checkmark & \checkmark & \checkmark & \checkmark & \checkmark & --         & \checkmark & 160 min  \\
\midrule
DART~\cite{huang2024dart}            & Implicit & RDA & --         & \checkmark & --         & --         & --         & \checkmark & \checkmark & 0.5 min  \\
Radar Fields~\cite{10.1145/3641519.3657510} & Implicit & RA & --   & \checkmark & --         & --         & --         & \checkmark & \checkmark & 1.6 min  \\
SpINR~\cite{takawale2025spinr}       & Implicit & RA  & --         & \checkmark & --         & --         & --         & \checkmark & \checkmark & --       \\
RadarSplat~\cite{kung2025radarsplat} & 3DGS     & RA  & --         & \checkmark & --         & --         & --         & \checkmark & \checkmark & 20 min   \\
\midrule
\textbf{Ours}                        & Points   & ADC & \checkmark & \checkmark & \checkmark & \checkmark & \checkmark & \checkmark & \checkmark & 2.3 min  \\
\bottomrule
\end{tabular}
}
\end{table}

Table~\ref{tab:comparison} summarizes the landscape of radar rendering methods along axes relevant to novel view synthesis. We organize related work into three categories corresponding to the approaches introduced in Sec.~\ref{sec:intro}.

\textbf{Physics-based inverse rendering for radar.}
Physics-based radar simulation has long employed shooting-and-bouncing rays (SBR) to produce path geometry and attenuation for large scenes and MIMO arrays~\cite{dudek2010millimeter}. These pipelines are forward-only: scene and sensor parameters are fixed, precluding calibration against real captures. Sionna RT~\cite{sionna} adds differentiability for antenna patterns and scattering coefficients but executes its path solver under gradient suspension, so vertex positions and surface normals receive no gradients; it also does not model wideband FMCW signals. Hofmann et al.~\cite{hofmann2025inverse} estimate stepped-frequency materials with per-vertex parameters but require known geometry, model only single-bounce propagation, and treat visibility as non-differentiable. InverTwin~\cite{chen2025invertwin} addresses multipath gradient pathologies via path-space differentiation but substitutes a Gaussian kernel surrogate for range processing and validates on synthetic data only. mmIR~\cite{mmIR2026} is the closest prior work: it places the full MC ray tracing pipeline inside an automatic differentiation graph, enabling end-to-end gradient flow from raw ADC through multi-bounce paths to per-vertex ITU materials, normals, and antenna patterns on real data. However, mmIR processes 24k surface interactions per iteration and requires ${\sim}$20 minutes for 500 iterations; it was designed for training-view rendering and virtual-aperture densification, not novel view synthesis. Our method replaces MC path tracing with deterministic hemisphere integral evaluation on 20k oriented points and complex range-profile splatting via a precomputed PSF, training all 8 viewpoints jointly in ${\sim}$2.3 minutes per scene (${\sim}$70$\times$ faster than mmIR's per-viewpoint training cost) while preserving the same ITU material physics and enabling NVS.

\textbf{Neural fields for radar novel view synthesis.}
Neural radiance fields (NeRF)~\cite{mildenhall2021nerf} have been adapted to radar across several modalities. DART~\cite{huang2024dart} trains a NeRF on range--Doppler images but discards phase entirely (magnitude-only supervision), cannot model FMCW waveforms, and does not expose material parameters. Radar Fields~\cite{10.1145/3641519.3657510} synthesizes power-only frequency-space responses via a neural field without physically interpretable scene parameters. SpINR~\cite{takawale2025spinr} models FMCW occupancy and Doppler recovery with an implicit neural representation. RF-4D~\cite{zhang2025rf4d} extends neural fields to dynamic scenes but remains sensor-specific. NeuRadar~\cite{rafidashti2025neuradar} and NeWRF~\cite{lu2024newrf} fuse radar with RGB or LiDAR modalities but do not render raw complex signals. All neural field methods share a structural limitation: they render post-processed images (RA maps, range--Doppler), requiring the network to implicitly learn FFT windowing artifacts (sidelobes, spectral leakage) and sensor-specific antenna patterns inside opaque features. None output complex-valued range profiles, so downstream tasks requiring custom signal processing (gesture tracking, vital-sign extraction, arbitrary beamforming) cannot be served without retraining. Our method renders complex range profiles that decouple scene physics from signal processing: identical FFT pipelines convert the output to any radar product.

\textbf{Gaussian splatting for radar.}
3D Gaussian splatting (3DGS)~\cite{3dgs} achieves real-time NVS in vision by rasterizing anisotropic 3D Gaussians into images. RadarSplat~\cite{kung2025radarsplat} adapts this primitive for RA rendering, learning per-Gaussian opaque features that predict radar intensity; it achieves fast inference but provides no physics-based material model, no complex output, and no mechanism to disentangle scene properties from sensor-specific artifacts. The underlying issue is that 3D Gaussians were designed for optical radiance fields, where each primitive emits view-dependent RGB via spherical harmonics. Porting this primitive to radar requires it to emit complex phasors at physically correct carrier phase and amplitude for each TX--RX pair, a fundamentally different output space that spherical harmonics do not naturally parameterize. Jiang et al.~\cite{jiang2025surfels} demonstrate the analogous point in optics: grounding the graphics primitive in the domain physics (Gaussian surfels with adapted radiosity) outperforms generic radiance field representations for relighting and geometry reconstruction. We make the same argument for radar. Our method uses oriented 3D points because they align directly with the radar hemisphere integral: each point stores per-surface ITU material parameters, evaluates a physically-based BSDF at the exact TX--RX geometry, and contributes a complex phasor to range-profile bins via Hann-windowed PSF splatting. This primitive choice yields physically interpretable materials, complex-valued output, and a factorized rendering pipeline amenable to fused CUDA acceleration (Table~\ref{tab:comparison}).
